\documentclass{article}
\usepackage{graphicx}
\usepackage{amsmath}
\usepackage{booktabs}
\usepackage{multirow}
\usepackage{array}
\usepackage{geometry}
\geometry{margin=1in}

\title{ORIE 5135 Project}
\author{Peter Ye}
\date{May 2026}

\begin{document}
\maketitle

\section{Question 1a: Natural ILP Formulation}

\textbf{Sets and Parameters:}
\begin{itemize}
    \item Set of exams: $V = \{1, 2, \dots, m\}$
    \item Set of available time slots: $T = \{1, 2, \dots, n\}$
    \item Number of proctors required by exam $i$: $r_i$
    \item Maximum proctor capacity per time slot: $R$
    \item Conflict edge set: $(i, i') \in E$ indicates exams $i$ and $i'$ share at least one student
\end{itemize}

\textbf{Decision Variables:}
\begin{itemize}
    \item $x_{ij} \in \{0,1\}$: equals 1 if exam $i$ is assigned to time slot $j$, 0 otherwise
    \item $y_j \in \{0,1\}$: equals 1 if time slot $j$ is used, 0 otherwise
\end{itemize}

\textbf{Model:}
\begin{align}
\min \quad & z = \sum_{j \in T} y_j \\
\text{s.t.} \quad & \sum_{j \in T} x_{ij} = 1, \quad \forall\, i \in V \\
& \sum_{i \in V} r_i\, x_{ij} \leq R\, y_j, \quad \forall\, j \in T \\
& x_{ij} + x_{i'j} \leq y_j, \quad \forall\, (i,i') \in E,\; \forall\, j \in T \\
& y_j \geq y_{j+1}, \quad \forall\, j = 1, \dots, n-1 \\
& x_{ij},\, y_j \in \{0,1\}, \quad \forall\, i \in V,\; j \in T
\end{align}

The last constraint is a symmetry-breaking inequality. Since the time slots are
interchangeable, any feasible schedule can be relabeled so that used slots come
before unused ones; this constraint enforces that canonical labeling, removing a
large amount of solution symmetry without changing the optimal objective value.

\section{Question 1b: Clique-Strengthened ILP Formulation}

The sets, parameters, decision variables, and all other constraints (including the
symmetry-breaking inequality) remain the same as in Formulation 1a. The pairwise
conflict constraints are replaced by clique inequalities.

Let $\mathcal{L}$ denote the set of all maximal cliques in the conflict graph $G$.

\textbf{Clique Inequality:}
\[
\sum_{i \in C} x_{ij} \leq y_j, \quad \forall\, C \in \mathcal{L},\; \forall\, j \in T
\]

\textbf{Validity and dominance.} Any two exams in a clique $C$ pairwise share at least
one student, so at most one exam from $C$ can be assigned to time slot $j$; if any
exam in $C$ is in slot $j$ then that slot is used, which gives $y_j = 1$. Hence
$\sum_{i \in C} x_{ij} \leq y_j$. Each such clique inequality dominates the
$\binom{|C|}{2}$ pairwise edge inequalities $x_{ij} + x_{i'j} \leq y_j$ for
$i, i' \in C$: summing the pairwise inequalities only yields
$\sum_{i \in C} x_{ij} \leq \tfrac{|C|}{2}\, y_j$, a strictly weaker bound for
$|C| \geq 3$. Replacing pairwise edge constraints with clique inequalities therefore
gives an LP relaxation that is at least as tight, and strictly tighter whenever the
LP solution would otherwise place fractional mass on three or more vertices of a
clique inside the same slot.

\section{Question 2a: Extended Formulation}

Let $\varphi$ denote the collection of all feasible groupings. A subset $S \subseteq \{1, 2, \dots, m\}$ belongs to $\varphi$ if and only if:
\begin{enumerate}
    \item \textbf{Independent set:} No two exams in $S$ share a student (i.e., $S$ is a stable set in $G$).
    \item \textbf{Capacity:} $\displaystyle\sum_{i \in S} r_i \leq R$.
\end{enumerate}

\textbf{Decision Variable:}
\begin{itemize}
    \item $\lambda_S \in \{0,1\}$: equals 1 if grouping $S$ is selected, 0 otherwise
\end{itemize}

\textbf{Model:}
\begin{align}
\min \quad & z = \sum_{S \in \varphi} \lambda_S \\
\text{s.t.} \quad & \sum_{\substack{S \in \varphi \\ i \in S}} \lambda_S = 1, \quad \forall\, i \in \{1, 2, \dots, m\} \\
& \lambda_S \in \{0,1\}, \quad \forall\, S \in \varphi
\end{align}

This is a set-partitioning model with one variable per feasible grouping, which is
exactly what we implement and use as the baseline in Question 3. We assume that the
number of available time slots $n$ is large enough to accommodate the optimal
schedule, so no explicit upper bound on the number of selected groupings is needed
(this holds for every instance in the two datasets, where $n = m$).

\section{Question 2b: Set Covering vs.\ Set Partitioning}

In the set partitioning formulation (used in Question 2a), each exam must belong to
\emph{exactly one} selected grouping ($= 1$). This requires enumerating \emph{all}
feasible groupings, including non-maximal ones, because restricting to only maximal
groupings may make it impossible to partition the exams such that each appears in
exactly one group.

An alternative is the set covering formulation, where each exam must belong to
\emph{at least one} selected grouping ($\geq 1$). Since the objective minimizes the
number of selected groupings, no optimal solution will redundantly cover an exam
in multiple groupings.

The set covering formulation allows us to restrict the variables to only
\textbf{maximal} feasible groupings --- groupings to which no additional exam can be
added without violating either the proctor capacity or a conflict constraint. This
significantly reduces the number of variables.

Despite this difference, the two formulations are equivalent: any solution of the
set partitioning model can be converted to a solution of the set covering model with
the same objective value (by expanding each non-maximal grouping to a maximal one),
and any solution of the set covering model can be converted back to a set
partitioning solution with the same value, by shrinking each maximal grouping to a
subset that ensures every exam is contained in exactly one of them. Their LP
relaxation values also coincide.

In practice, the set covering formulation with maximal groupings is preferable because:
\begin{enumerate}
    \item The number of variables is smaller (though still exponential in the worst case).
    \item The LP relaxation values are identical.
    \item LPs with inequality constraints ($\geq$) are generally easier for solvers
          than LPs with equality constraints ($=$).
\end{enumerate}

The set covering formulation replaces the equality constraints with inequalities:

\begin{align}
\min \quad & z = \sum_{S \in \mathcal{S}} \lambda_S \\
\text{s.t.} \quad & \sum_{\substack{S \in \mathcal{S} \\ i \in S}} \lambda_S \geq 1, \quad \forall\, i \in \{1, 2, \dots, m\} \\
& \lambda_S \in \{0,1\}, \quad \forall\, S \in \mathcal{S}
\end{align}

where $\mathcal{S}$ denotes the collection of all \textbf{maximal} feasible groupings. A feasible grouping $S$ is maximal if no additional exam can be added to $S$ without violating either the proctor capacity constraint or the independence requirement.

We did not implement the set covering version, but compared with set partitioning we
would expect it to be faster on both ends of the pipeline. Enumeration would produce
fewer columns, since only maximal groupings need to be stored rather than every
smaller subset visited along the way. Each LP solve would also be slightly cheaper,
because there are fewer variables and inequality constraints are usually easier for
the solver to handle than equalities. The root LP value stays the same, so the bound
quality is unchanged.

\section{Question 3: Computational Experiments}

All experiments were run with Gurobi (via \texttt{gurobipy}) on a single machine. The
implementation reads each instance directly from its JSON file and builds the three
formulations described in Sections 1 and 2. For every formulation, we record (i) the LP
relaxation value at the root, (ii) the number of branch-and-bound nodes explored, and
(iii) the wall-clock solve time --- as reported by Gurobi for Formulations 1a and 1b,
and as total wall time including the enumeration of feasible groupings for
Formulation~2. Each instance has $m = n$ and $R = 25$, with conflict-graph density
close to $0.5$.

\subsection{Question 3(a): Dataset 1 (time limit 120\,s)}

Dataset 1 contains ten instances of growing size, ranging from $m = 35$
($\lvert E \rvert = 292$) to $m = 75$ ($\lvert E \rvert \approx 1400$). Table~\ref{tab:ds1}
reports the results for the three formulations on every instance.

\begin{table}[h]
\centering
\footnotesize
\caption{Dataset 1 results. ``Opt?'' indicates whether Gurobi proved optimality within the
120\,s time limit. ``LP'' is the value of the LP relaxation at the root.
For Formulation~2, ``Time'' includes the time spent enumerating feasible groupings.}
\label{tab:ds1}
\begin{tabular}{llrrrrcrr}
\toprule
Instance & ($m$, $\lvert E\rvert$) & Form. & LP & Obj. & Opt? & Nodes & Time (s) \\
\midrule
\multirow{3}{*}{01} & \multirow{3}{*}{(35, 292)}
   & 1a Natural  &  6.32 &  8 & Y & 1 &   1.20 \\
 & & 1b Clique   &  6.32 &  8 & Y & 1 &   2.35 \\
 & & 2 Extended  &  7.18 &  8 & Y & 1 &   0.04 \\
\midrule
\multirow{3}{*}{02} & \multirow{3}{*}{(38, 345)}
   & 1a Natural  &  7.32 &  9 & Y &  943 &   1.89 \\
 & & 1b Clique   &  7.32 &  9 & Y &  791 &   7.94 \\
 & & 2 Extended  &  7.92 &  9 & Y &    1 &   0.08 \\
\midrule
\multirow{3}{*}{03} & \multirow{3}{*}{(40, 369)}
   & 1a Natural  &  7.88 &  9 & Y &  209 &   2.69 \\
 & & 1b Clique   &  7.88 &  9 & Y &  153 &   3.24 \\
 & & 2 Extended  &  8.12 &  9 & Y &    1 &   0.10 \\
\midrule
\multirow{3}{*}{04} & \multirow{3}{*}{(60, 886)}
   & 1a Natural  & 12.84 & 13 & Y & 207\,271 &  99.36 \\
 & & 1b Clique   & 12.84 & 13 & N & 9\,104   & 120.05 \\
 & & 2 Extended  & 12.84 & 13 & Y & 1        &   0.61 \\
\midrule
\multirow{3}{*}{05} & \multirow{3}{*}{(63, 947)}
   & 1a Natural  & 12.40 & 13 & Y & 1 &  10.89 \\
 & & 1b Clique   & 12.40 & 13 & Y & 1 &  37.28 \\
 & & 2 Extended  & 12.40 & 13 & Y & 1 &   0.58 \\
\midrule
\multirow{3}{*}{06} & \multirow{3}{*}{(65, 1057)}
   & 1a Natural  & 12.88 & 14 & N & 32\,091 & 120.04 \\
 & & 1b Clique   & 12.88 & 14 & N & 3\,781  & 120.03 \\
 & & 2 Extended  & 12.88 & 13 & Y & 413     &   2.44 \\
\midrule
\multirow{3}{*}{07} & \multirow{3}{*}{(68, 1109)}
   & 1a Natural  & 12.68 & 13 & Y & 6\,516 &  21.26 \\
 & & 1b Clique   & 12.68 & 13 & Y & 1\,524 & 101.51 \\
 & & 2 Extended  & 12.68 & 13 & Y & 1      &   1.76 \\
\midrule
\multirow{3}{*}{08} & \multirow{3}{*}{(75, 1380)}
   & 1a Natural  & 15.36 & 16 & Y & 1 &  20.02 \\
 & & 1b Clique   & 15.36 & 16 & Y & 1 & 115.48 \\
 & & 2 Extended  & 15.36 & 16 & Y & 1 &   0.77 \\
\midrule
\multirow{3}{*}{09} & \multirow{3}{*}{(75, 1372)}
   & 1a Natural  & 14.16 & 15 & Y & 1 &  15.44 \\
 & & 1b Clique   & 14.16 & 75 & N & 1 & 120.02 \\
 & & 2 Extended  & 14.16 & 15 & Y & 1 &   1.13 \\
\midrule
\multirow{3}{*}{10} & \multirow{3}{*}{(75, 1423)}
   & 1a Natural  & 14.68 & 16 & N & 29\,094 & 120.01 \\
 & & 1b Clique   & 14.68 & 73 & N & 1       & 120.04 \\
 & & 2 Extended  & 14.68 & 15 & Y & 1       &   0.89 \\
\bottomrule
\end{tabular}
\end{table}

\paragraph{(i) LP relaxation values.}
The natural formulation (1a) and the clique-strengthened formulation (1b) produce
\emph{identical} root LP values on every instance of Dataset~1 (they differ only at
the last decimal due to floating-point noise). At first glance this is surprising,
since the clique inequality strictly dominates the corresponding pairwise edge
inequalities. The explanation has two parts. First, on these graphs with edge density
$\sim\!0.5$, the LP optimum is largely determined by the proctor capacity bound
$\sum_i r_i / R$ together with the symmetry-breaking constraints $y_j \geq y_{j+1}$;
once the LP saturates these constraints, the additional clique inequalities are not
active. Second, Gurobi's presolve very likely generates clique cuts from the pairwise
edge constraints automatically, so writing the cliques explicitly does not add new
information at the root LP. Formulation~2 is provably at least as tight as
Formulations 1a and 1b, and we observe \emph{strict} improvement only on the
smallest, sparsest instances (01--03), where the conflict-graph structure dominates
over the capacity bound; from instance~04 onward the capacity bound becomes binding
and all three formulations report the same LP value.

\paragraph{(ii) Branch-and-bound nodes vs.\ solve time.}
Equal LP relaxation values do \emph{not} imply equal branch-and-bound effort. The clearest
example is instance~04: all three formulations have LP value $12.84$, but Formulation~2
closes the gap at the root (1 node, 0.6\,s), Formulation~1b explores
$\sim\!9\,000$ nodes and times out, and Formulation~1a explores
$\sim\!2\!\times\!10^5$ nodes (99\,s). The same pattern appears on instances 02, 06,
and 07. Two factors drive this. (a) The model size differs dramatically: 1b adds one
constraint per (maximal clique, time slot) pair, which on dense graphs yields many
thousands of additional rows, making each node LP slower to solve and the
branching/heuristic stage less effective at finding incumbents. (b) The extended
formulation has very few variables relative to its information content (each
$\lambda_S$ encodes an entire feasible slot), so Gurobi's primal heuristics typically
recover an integer-optimal incumbent immediately at the root. Thus our data suggests
that LP bound quality is necessary but not sufficient for fast B\&B; the structural
compactness of the formulation and the difficulty of each individual node LP matter
just as much. A striking failure mode is visible on instances 09 and 10 of Formulation~1b,
which return objective values of 75 and 73 respectively---near the trivial
``one exam per slot'' upper bound of $m = 75$---meaning that within the 120\,s limit
Gurobi's heuristics could not find a tight feasible solution under this formulation.

\paragraph{(iii) Scaling with instance size.}
The trend across the ten instances is clear. Formulation~2 (extended) scales the best
\emph{in solve time}: nine of the ten instances are closed with a single B\&B node,
and total time stays below $2.5$\,s. The bottleneck of Formulation~2 is purely
combinatorial enumeration of the feasible groupings, which is still tractable for
$m \leq 75$ at this density but grows exponentially. Formulations 1a and 1b both
suffer rapidly increasing node counts as $m$ and $|E|$ grow: the natural formulation
times out on instances 06 and 10 and the clique formulation times out on instances
04, 06, 09, and 10. The clique formulation typically explores fewer nodes than the
natural formulation when both finish (e.g., 1\,524 vs.\ 6\,516 on instance~07), but
its per-node cost is much higher because of the larger constraint matrix, so its
wall-clock time is consistently worse on this dataset. In summary, growth rates differ
substantially: Formulation~2 grows in \emph{enumeration} cost, while Formulations 1a
and 1b grow in \emph{B\&B} cost, with 1b having the worst wall-clock growth in our
experiments.

\subsection{Question 3(b): Dataset 2 (time limit 600\,s)}

Dataset 2 contains nine larger instances with $m \in \{80, \ldots, 90\}$ and
$\lvert E \rvert$ up to $\sim\!2000$. Table~\ref{tab:ds2} reports the results for
Formulations 1a and 1b. Formulation~2 is omitted because the enumeration of feasible
groupings is no longer practical at this scale (see discussion below).

\begin{table}[h]
\centering
\footnotesize
\caption{Dataset 2 results (time limit 600\,s).}
\label{tab:ds2}
\begin{tabular}{llrrrrcrr}
\toprule
Instance & ($m$, $\lvert E\rvert$) & Form. & LP & Obj. & Opt? & Nodes & Time (s) \\
\midrule
\multirow{2}{*}{01} & \multirow{2}{*}{(80, 1584)}
   & 1a Natural & 16.32 & 17 & Y & 1 &  22.26 \\
 & & 1b Clique  & 16.32 & 17 & Y & 1 & 205.09 \\
\midrule
\multirow{2}{*}{02} & \multirow{2}{*}{(82, 1682)}
   & 1a Natural & 16.08 & 17 & Y & 1 &  15.55 \\
 & & 1b Clique  & 16.08 & 17 & Y & 1 & 347.32 \\
\midrule
\multirow{2}{*}{03} & \multirow{2}{*}{(83, 1726)}
   & 1a Natural & 16.28 & 17 & Y & 1 &  32.54 \\
 & & 1b Clique  & 16.28 & 17 & Y & 1 & 355.03 \\
\midrule
\multirow{2}{*}{04} & \multirow{2}{*}{(87, 1843)}
   & 1a Natural & 16.76 & 17 & Y & 9\,995 &  75.71 \\
 & & 1b Clique  & 16.76 & 17 & Y & 1      & 556.95 \\
\midrule
\multirow{2}{*}{05} & \multirow{2}{*}{(87, 1867)}
   & 1a Natural & 16.60 & 17 & Y & 2\,962 &  64.76 \\
 & & 1b Clique  & 16.60 & 81 & N & 1      & 600.09 \\
\midrule
\multirow{2}{*}{06} & \multirow{2}{*}{(87, 1821)}
   & 1a Natural & 17.76 & 18 & Y & 3\,434 &  51.99 \\
 & & 1b Clique  & 17.76 & 18 & N & 216    & 600.14 \\
\midrule
\multirow{2}{*}{07} & \multirow{2}{*}{(82, 1660)}
   & 1a Natural & 16.88 & 17 & Y & 20\,031 &  97.49 \\
 & & 1b Clique  & 16.88 & 18 & N & 2\,347  & 600.08 \\
\midrule
\multirow{2}{*}{08} & \multirow{2}{*}{(85, 1805)}
   & 1a Natural & 16.84 & 17 & Y & 37\,245 & 214.19 \\
 & & 1b Clique  & 16.84 & 18 & N & 1       & 600.06 \\
\midrule
\multirow{2}{*}{09} & \multirow{2}{*}{(90, 2032)}
   & 1a Natural & 18.00 & 19 & N & 153\,156 & 600.03 \\
 & & 1b Clique  & 18.00 & 20 & N & 1        & 600.12 \\
\bottomrule
\end{tabular}
\end{table}

\paragraph{Comparison and computational limits.}
The qualitative pattern observed on Dataset~1 is amplified here. The two LP
relaxations remain identical on every instance, again because the capacity bound and
the symmetry-breaking constraints dominate the structure of the LP optimum and
Gurobi's presolve already exploits clique structure internally. In terms of solving the
ILP, however, Formulation~1a clearly outperforms Formulation~1b: 1a closes eight of
the nine instances within the time limit (only the largest, instance~09, times out),
whereas 1b proves optimality on only \emph{four} of the nine. On the remaining five
instances, 1b often returns notably weaker incumbents (e.g., 81 on instance~05, 20 on
instance~09 against 1a's 19), reflecting that the heavy constraint matrix slows both
LP solves at each node and primal heuristics, so the solver cannot find tight integer
solutions in time.
This is the same trade-off seen on Dataset~1, scaled up: clique strengthening that
should in principle help branching becomes counterproductive once the number of maximal
cliques explodes on dense graphs of this size. Even the natural formulation 1a starts
to scale poorly: solve time grows from $\sim\!20$\,s at $m = 80$ to over $200$\,s at
$m = 85$, with $1.5\!\times\!10^5$ nodes explored on the largest instance before
timing out. Both formulations are therefore approaching their practical limit at
$m \approx 90$ on graphs of edge density $\sim\!0.5$.

\paragraph{Why the extended formulation is not applicable.}
The extended formulation has one variable per feasible grouping, where a grouping is
both an independent set in $G$ and a knapsack-feasible subset under the proctor budget.
On Dataset~1 (up to $m = 75$) the recursive enumeration in our implementation is still
feasible. On Dataset~2 the situation changes qualitatively: with $m \approx 90$ and
edge density $\sim\!0.5$, the conflict graph contains many large independent sets, and
the proctor capacity $R = 25$ is not tight enough to prune the search aggressively.
The number of feasible groupings grows exponentially in $m$ on graphs of this density,
so on these instances it is already astronomically large; even storing the complete
variable set in memory is impractical, let alone passing it to Gurobi. In other words, the extended formulation trades a polynomial-size model
(Formulations 1a/1b) for an exponentially large one whose strength comes precisely
from that exhaustive enumeration; once enumeration is no longer feasible, the
formulation cannot be built explicitly, which is why Question~3(b) restricts the
comparison to the two compact formulations.

\end{document}